\documentclass[a4paper, 12pt]{article}

\usepackage{utils}

\renewcommand*{\today}{21 January 2026}

\begin{document}

\hotbox{Géometrie 2}{CM 2}{\today}

\begin{theoreme}{}{}
    \begin{enumerate}
        \item Soit $\F \subset \E$ tel que $\F \neq \emptyset$. Alors $\F$ est un sous-espace affine de $\E$ si et seulement si tout barycentre d'un système de points pondérés de
        $\F$ est encore un point de $\F$.

        \item Soit $X \in \E$ tel que $X \neq \emptyset$ alors $Aff(X) = A + Vect(\overrightarrow{AM} | M \in X)$ pour un $A \in X$ est l'ensemble
        des barycentres $bar(\mathcal{A})$ où $\mathcal{A}$ parcourt l'ensemble des points pondérés de $X$.
    \end{enumerate}
\end{theoreme}

\begin{demonstration}
    \begin{enumerate}
        \item 
        \begin{description}
            \item[Sens direct.] 
            Si $\F$ est un sous-espace affine de $\E$ et si $G = bar(\mathcal{A})$ où $\mathcal{A} = ((A_1, \alpha_1), \ldots, (A_n, \alpha_n))$ et $A_i \in \F$,
            Alors $\overrightarrow{A_1 G} = \dfrac{1}{\alpha} \sum_{i=1}^n \alpha_i \overrightarrow{A_1 A_i} = \dfrac{1}{\alpha} \sum_{i=2}^n \alpha_i \overrightarrow{A_1 A_i}$
            d'où $G = A_1 + \sum_{i=2}^n \dfrac{\alpha_i}{\alpha} \overrightarrow{A_1 A_i}$ mais $\overrightarrow{A_1 A_i} \in \overrightarrow{\F}$, $\dfrac{a_i}{\alpha} \in \R$
            et $\sum_{i=2}^n \dfrac{\alpha_i}{\alpha} \in \overrightarrow{\F}$ donc $G \in \F$ car il s'écrit comme $A_1 + \vec{v}$ avec $A_1 \in \F$ et $\vec{v} \in \overrightarrow{\F}$.
            \item[Sens réciproque.]
            On suppose que $\F$ est stable par la prise de barycentre et on veut montrer que $\F$ est un sous-espace affine.
            Comme $\F \neq \emptyset$, il existe $A \in \F$.

            On considère $\overrightarrow{V} = \{ \overrightarrow{AM} | M \in \F \}$. Montrons que $\overrightarrow{V}$ est un sous-espace vectoriel de $\overrightarrow{E}$.
            On peut écrire $\overrightarrow{AA} = \overrightarrow{0} \in \overrightarrow{V}$ donc $\overrightarrow{V} \neq \emptyset$.
            Soient $\overrightarrow{AM}, \overrightarrow{AN} \in \overrightarrow{V}$ et $\lambda\in \R$.
            On a 
            \begin{align*}
                \overrightarrow{AM} + \lambda \overrightarrow{AN} = \overrightarrow{AP} &\iff \overrightarrow{AP} = \overrightarrow{AP} + \overrightarrow{PM} + \lambda (\overrightarrow{AP} + \overrightarrow{PN}) \\
                &\iff \vec{0} = \overrightarrow{PM} + \lambda \overrightarrow{PN} - \lambda \overrightarrow{PA} \\
                &\iff P = bar((M, 1), (N, \lambda), (A, -\lambda))
            \end{align*}
            où $P = bar((M, \lambda), (N, 1 - \lambda), (A, 0))$.
            Par hypothèse, $P \in \F$ donc $\overrightarrow{AP} \in \overrightarrow{V}$.
            Ainsi, $\overrightarrow{V}$ est un sous-espace vectoriel de $\overrightarrow{E}$.
            $\F = A + Vect(\overrightarrow{AM} | M \in \F)$ est bien un sous-espace affine avec $\overrightarrow{\F} = \overrightarrow{V}$.
        \end{description}
    \end{enumerate}
\end{demonstration}

\begin{definition}[(Repère affine)]
    Soit $\E$ un espace affine de dimension $n$ et de direction $\vec{\E}$. Un \textbf{repère affine} de $\E$ est une famille $(A_0, A_1, \ldots, A_n)$ de points de $\E$ tels que
    $(\overrightarrow{A_0 A_1}, \ldots, \overrightarrow{A_0 A_n})$ soit une base de $\vec{\E}$.

    Ainsi pour tout $M \in \E$, il existe des réels uniques $(\alpha_1, \ldots, \alpha_n)$ tels que
    $$
    \overrightarrow{A_0 M} = \sum_{i=1}^n \alpha_i \overrightarrow{A_0 A_i}
    $$
    Les réels $\alpha_i$ sont appelés les \textbf{coordonnées barycentriques} de $M$ dans le repère affine $(A_0, A_1, \ldots, A_n)$.
\end{definition}

\begin{methode}
    \textbf{Passage entre coordonnées barycentriques et coordonnées dans un repère affine}
    \begin{align*}
        \sum_{i=0}^n \alpha_i \overrightarrow{M A_i} = \vec{0} &\iff \sum_{i=0}^n \alpha_i (\overrightarrow{MA_0} + \overrightarrow{A_0 A_i}) = \vec{0} \\
        &\iff \alpha_i \overrightarrow{M A_0} + \sum_{i=0}^n \alpha_i \overrightarrow{A_0 A_i} = \vec{0} \\
        &\iff \sum_{i=0}^n \alpha_i \overrightarrow{A_0 A_i} = (\sum_{i=0}^n \alpha_i) \overrightarrow{A_0 M} \\
        &\iff \overrightarrow{A_0 M} = \sum_{i=1}^n \alpha_i \overrightarrow{A_0 A_i}
    \end{align*}
\end{methode}

\end{document}